% Сборка: xelatex docs/dsl_full_report.tex (дважды для оглавления)
% Требуются: XeLaTeX, polyglossia, fontspec, шрифты Times New Roman / Arial / Consolas
% (поставляются с Windows). Для Linux/macOS можно заменить на DejaVu Serif/Sans/Mono.

\documentclass[12pt,a4paper]{article}

\usepackage{fontspec}
\usepackage{polyglossia}
\setdefaultlanguage{russian}
\setotherlanguage{english}

\setmainfont{Times New Roman}
\setsansfont{Arial}
\setmonofont{Consolas}
\newfontfamily\cyrillicfonttt{Consolas}

\usepackage[a4paper,margin=2.5cm]{geometry}
\usepackage{amsmath}
\usepackage{fancyvrb}
\usepackage{xcolor}
\usepackage{enumitem}
\usepackage{titlesec}
\usepackage{hyperref}
\hypersetup{
  colorlinks=true,
  linkcolor=black,
  urlcolor=blue,
  citecolor=blue,
  pdftitle={DSL подсистема DiGr. Полный отчёт},
  pdfauthor={Дрекалов Н.С., Соколов А.Н., Зеленов Н.М.}
}

\definecolor{dslkw}{RGB}{0,90,156}
\definecolor{boxbg}{RGB}{245,245,250}

\setlist[itemize]{leftmargin=1.6em,topsep=2pt,itemsep=2pt}
\setlist[enumerate]{leftmargin=1.8em,topsep=2pt,itemsep=2pt}

\titlespacing*{\section}{0pt}{1.4em}{0.6em}
\titlespacing*{\subsection}{0pt}{1.0em}{0.4em}
\titlespacing*{\subsubsection}{0pt}{0.8em}{0.3em}

\newcommand{\uml}[1]{\textit{См. диаграмму \texttt{docs/uml/#1}.}}
\newcommand{\code}[1]{\texttt{#1}}

\title{DSL подсистема DiGr.\\ Полный технический отчёт}
\author{Дрекалов Н.\,С. \and Соколов А.\,Н. \and Зеленов Н.\,М.}
\date{Апрель 2026}

\begin{document}
\maketitle
\tableofcontents
\newpage

% =============================================================
\section{Введение}

\subsection{Контекст и постановка задачи}

DiGr (Document Graph) есть инструментарий для построения и анализа AST
произвольных текстовых документов: TeX, plain text, потенциально других
форматов. Парсер документа конфигурируется YAML-файлами в
\code{config/formats/} и строит дерево \code{AstDocument} из узлов
\code{AstNode} (см. модуль \code{src/document\_ast/}).

Над готовым AST требуется выполнять структурные запросы трёх классов:

\begin{enumerate}
  \item \textbf{FIND}: выбор узлов по типу сущности и предикату.
  \item \textbf{CONTEXT}: поиск минимальных окон базовых единиц
        текста, содержащих заданные паттерны.
  \item \textbf{DISTANCE}: измерение расстояний между парами
        узлов в выбранных единицах (символы или экземпляры заданной
        сущности).
\end{enumerate}

Для решения этой задачи в проекте был спроектирован и реализован
DSL (domain-specific language) с собственным синтаксисом, типобезопасным
AST запроса и акторным конвейером исполнения. Настоящий отчёт
последовательно описывает: что было реализовано, в каком порядке, как
именно и зачем, а также демонстрирует работу системы на реальных
запросах из \code{examples.md}.

\uml{component\_overview.puml}

\subsection{Структура отчёта}

Отчёт разбит на четыре содержательные части:

\begin{itemize}
  \item раздел~2: хронология разработки с обоснованием
        порядка реализации.
  \item раздел~3: архитектура DSL и реализационные детали
        каждого слоя.
  \item раздел~4: акторный конвейер парсинга и исполнения.
  \item раздел~5: четыре примера запросов из
        \code{examples.md} вместе с фактическим выводом CLI на
        документе \code{GA\_1\_2025.tex}.
\end{itemize}

% =============================================================
\section{Что и в каком порядке было реализовано}

\subsection{Хронология}

Порядок реализации диктуется направлением зависимостей: сначала
\code{actor}, на нём поднимается \code{document\_ast}, а уже на этой
базе строится \code{dsl}. Никакой нижний слой не знает о вышестоящем.
Хронологически работа шла так:

\begin{enumerate}
  \item \textbf{Actor framework} (\code{src/actor/}). Базовый runtime
        с типизированными акторами (\code{Actor[State, QueueItem,
        Message]}), драйвером \code{ProceedableActorDriver}, очередями
        \code{Mailbox} и хэндлами \code{ActorHandle}. Без него нельзя
        строить конвейеры с явными FSM-состояниями.
        \uml{actor\_architecture.puml}
  \item \textbf{Document AST} (\code{src/document\_ast/}). Модель
        документа (\code{AstDocument}, \code{AstNode}) и акторный
        парсер, использующий конфигурации форматов из
        \code{config/formats/}. Сегментеры (\code{plain\_text},
        \code{latex\_content\_scope}, \code{latex\_semantic\_block},
        \code{latex\_definition}, \dots) задают, как разбивать текст
        на узлы.
  \item \textbf{DSL core} (\code{src/dsl/}). Модель запроса, лексер,
        recursive-descent парсер, исполнитель. Реализован в едином
        коммите после готовности фреймворка акторов и AST.
  \item \textbf{Поддержка TeX-формата}: \code{config/formats/tex.yaml}
        и набор LaTeX-сегментеров. Это позволило применять DSL к
        реальным учебным TeX-документам.
  \item \textbf{Запросы DISTANCE}: измерение расстояний между
        парами узлов с поддержкой различных единиц измерения и
        статистики по выборке.
  \item \textbf{Сущность \code{definition}}: распознавание
        именованных определений (макросов вида \code{\textbackslash odmkey})
        как отдельных узлов AST.
  \item \textbf{Фикс CONTEXT-запроса}: корректная обработка
        вырожденных и минимальных окон.
\end{enumerate}

\subsection{Почему именно такой порядок}

\begin{itemize}
  \item \textbf{Снизу вверх по зависимостям.} Если бы DSL писался
        раньше акторов, его пришлось бы переписывать под асинхронную
        модель; начав с фреймворка акторов, авторы получили готовую
        абстракцию для конвейеров.
  \item \textbf{Сначала общий каркас, потом форматы.} TeX-конфиг
        написан после ядра DSL, потому что DSL не должен ничего знать
        о конкретных форматах: он работает с уже построенным
        \code{AstDocument}.
  \item \textbf{Запросы поэтапно.} Сначала \code{FIND} как минимальный
        полезный запрос, затем \code{CONTEXT} как развитие идеи окна,
        затем \code{DISTANCE} как самостоятельный аналитический
        примитив. Каждый шаг расширяет грамматику и исполнитель, но
        не ломает старое.
\end{itemize}

% =============================================================
\section{Архитектура и реализация}

\subsection{Слои подсистемы}

В \code{src/dsl/} выделены шесть семантических слоёв:

\begin{itemize}
  \item \code{api/}: публичный фасад \code{ActorDslEngine}
        (\code{api/engine.py}), \code{ActorDslParser}, \code{ActorDslExecutor}.
  \item \code{model/}: типобезопасный AST запроса
        (\code{model/query\_ast.py}).
  \item \code{parsing/}: лексер, токены, recursive-descent
        парсер, типизированные ошибки.
  \item \code{execution/}: индекс документа, вычислитель
        предикатов, калькулятор расстояний, валидатор запроса,
        структуры результата.
  \item \code{actors/parsing/}: акторы конвейера разбора.
  \item \code{actors/execution/}: акторы конвейера исполнения.
\end{itemize}

\uml{dsl\_architecture.puml}

\subsection{Модель запроса (AST)}

Файл: \code{src/dsl/model/query\_ast.py}.

Корневой sum-тип:
\[
\code{DslQuery} = \code{ContextQuery}\ |\ \code{FindQuery}\ |\ \code{DistanceQuery}.
\]

Sum-тип для предикатов \code{WHERE}:
\[
\code{Expression} = \code{ComparisonExpression}\ |\ \code{FunctionExpression}
\ |\ \code{NotExpression}\ |\ \code{BinaryExpression}.
\]

Вспомогательные узлы: \code{Selector}, \code{Pattern}, \code{SpanSpec},
\code{WithinConstraint}, \code{FieldRef}, \code{CountConstraint},
\code{RegexLiteral}, \code{DistanceReturn}, \code{PairLimit}. Все они
наследуют миксин \code{Serializable} и поддерживают \code{to\_dict()}.

\textbf{Зачем такая модель:}

\begin{itemize}
  \item Дата-классы вместо словарей дают типобезопасность и
        автодополнение в IDE.
  \item Сериализация через \code{to\_dict()} нужна для логирования,
        отладки и тестов золотых снимков.
  \item Sum-типы делают разбор ветвлений в исполнителе исчерпывающим:
        пропуск варианта виден на этапе ревью.
\end{itemize}

\subsection{Парсинг: лексер и recursive descent}

\subsubsection{Лексер}

Файл: \code{src/dsl/parsing/lexer.py}. Класс \code{DslLexer} разбивает
текст на токены \code{DslToken(kind, value, position, line, column)}.
Виды токенов перечислены в \code{token\_kind.py}:

\begin{itemize}
  \item ключевые слова: \code{CONTEXT}, \code{FIND}, \code{DISTANCE},
        \code{FOR}, \code{TO}, \code{WITHIN}, \code{LIMIT\_PAIRS},
        \code{WHERE}, \code{RETURN}, \code{AND}, \code{OR}, \code{NOT};
  \item операторы сравнения: \code{\textasciitilde =},
        \code{!\textasciitilde =}, \code{=}, \code{!=}, \code{<},
        \code{>}, \code{<=}, \code{>=};
  \item литералы: идентификаторы, строки в двойных кавычках,
        regex-литералы, целые числа;
  \item пунктуация: скобки, запятые, точки.
\end{itemize}

\subsubsection{Парсер}

Файл: \code{src/dsl/parsing/recursive\_descent\_parser.py}. Класс
\code{DslTokenStreamParser} реализован методом рекурсивного спуска.
Грамматика небольшая, нет леворекурсивных правил, поэтому recursive
descent уместен и легко поддерживается.

Приоритеты булевых операций в \code{WHERE}:
\code{NOT} > \code{AND} > \code{OR}. Селекторы парсятся в
\code{Selector(entity, predicate)}, вызовы функций в
\code{FunctionExpression(name, args)}.

Ошибки агрегируются в \code{DslSyntaxError(message, offset, line,
column)}, благодаря чему пользователь видит точную позицию в исходнике.

\uml{dsl\_parser\_sequence.puml}

\subsubsection{Почему recursive descent}

\begin{itemize}
  \item Грамматика DSL компактна и не требует генератора парсеров.
  \item Ручной парсер легко расширять (как и было сделано при
        добавлении \code{DISTANCE}).
  \item Сообщения об ошибках формируются под конкретное место,
        а UX в итоге лучше, чем у автогенерированного парсера.
\end{itemize}

\subsection{Исполнение запроса}

\subsubsection{DocumentIndex}

Файл: \code{src/dsl/execution/document\_index.py}. Создаётся один раз
на документ. При построении проходит AST и индексирует:

\begin{itemize}
  \item все узлы по типу сущности (\code{nodes\_of\_entity});
  \item отношения parent/child (для быстрого \code{children},
        \code{descendants}, \code{ancestors});
  \item диапазоны узлов внутри произвольного спана
        (\code{nodes\_within\_span}, \code{container\_nodes\_for\_span});
  \item соседей по позиции (\code{previous\_node}, \code{next\_node}).
\end{itemize}

Сложность построения \(O(n)\), где \(n\) есть количество узлов; запросы
по типу сущности выполняются за \(O(1)\) с амортизацией.

\subsubsection{PredicateEvaluator}

Файл: \code{src/dsl/execution/predicate\_evaluator.py}. Вычисляет
\code{WHERE} над текущим узлом. Поддерживает:

\begin{itemize}
  \item операторы сравнения над полями метаданных и спанов;
  \item regex-сопоставление через \code{\textasciitilde =} и
        \code{!\textasciitilde =};
  \item встроенные функции: \code{has\_child}, \code{has\_descendant},
        \code{has\_ancestor}, \code{contains}, \code{follows},
        \code{precedes}, \code{intersects}, \code{inside};
  \item сопоставление паттернов внутри текстовых окон.
\end{itemize}

\subsubsection{DistanceCalculator}

Файл: \code{src/dsl/execution/distance\_calculator.py}. Генерирует пары
узлов по двум селекторам, отбрасывает пары, не удовлетворяющие
ограничениям, и считает расстояние:

\begin{itemize}
  \item \code{distance(symbol)}: разница позиций в символах
        между ближайшими краями спанов;
  \item \code{distance(<entity>)}: количество узлов указанной
        сущности, лежащих между парой (используется бинпоиск
        \code{bisect\_left} по отсортированным позициям, что даёт
        \(O(\log n)\) на пару);
  \item кэш диапазонов по типам сущностей переиспользуется между
        парами одной выборки.
\end{itemize}

\subsubsection{QueryValidator и структуры результата}

\code{QueryValidator} перед запуском проверяет, что все упомянутые в
запросе сущности есть в документе. Тогда ошибка вылезает заранее,
а не в середине обхода.

Файл: \code{src/dsl/execution/query\_results.py} определяет:

\begin{itemize}
  \item \code{FindMatch}, \code{ContextWindowMatch},
        \code{DistancePairMatch}: элементы выборки;
  \item \code{FindQueryExecutionResult},
        \code{ContextQueryExecutionResult},
        \code{DistanceQueryExecutionResult}: обёртки с полями
        \code{count}, \code{returns}, \code{items} и опциональным
        \code{stats};
  \item \code{PatternMatchResult}: результат сопоставления
        паттерна внутри окна.
\end{itemize}

\uml{dsl\_executor\_sequence.puml}

% =============================================================
\section{Акторный конвейер}

\subsection{Общая схема}

И парсинг, и исполнение организованы единообразно: координатор
раздаёт работу воркерам, сборщик аккумулирует результат. Логическая
цепочка такова: \code{Coordinator}, далее \code{Workers}~$1{\dots}N$,
далее \code{Collector}.

Каждый актор имеет явный FSM, отражающий доменные фазы, а не общие
\code{idle/working}. Это упрощает диагностику и тестирование.

\uml{actor\_architecture.puml}

\subsection{Конвейер парсинга}

Файлы: \code{src/dsl/actors/parsing/}.

\begin{itemize}
  \item \code{DslLexerActor}: оборачивает \code{DslLexer},
        принимает \code{TokenizeDslRequest}, возвращает поток токенов.
  \item \code{DslQueryParserActor}: оборачивает
        \code{DslTokenStreamParser}; FSM проходит фазы
        \code{CLASSIFYING\_QUERY}, \code{PARSING\_FIND\_TARGET},
        \code{PARSING\_FIND\_CONSTRAINTS}, \code{PARSING\_CONTEXT\_SPAN},
        \code{PARSING\_CONTEXT\_PATTERNS},
        \code{PARSING\_CONTEXT\_CONSTRAINTS}.
  \item \code{DslParserCoordinatorActor}: оркеструет цепочку
        лексер, парсер, сборщик; FSM проходит состояния
        \code{IDLE}, \code{WAITING\_FOR\_TOKENS},
        \code{WAITING\_FOR\_QUERY\_AST}, \code{COMPLETED}.
  \item \code{DslParseResultCollectorActor}: собирает финальный
        AST либо ошибку.
\end{itemize}

\uml{dsl\_parser\_sequence.puml}

\subsection{Конвейер исполнения}

Файлы: \code{src/dsl/actors/execution/}.

\begin{itemize}
  \item \code{DslExecutionCoordinatorActor}: главный оркестратор;
        FSM проходит состояния \code{IDLE},
        \code{EVALUATING\_FIND\_CANDIDATES} либо
        \code{EVALUATING\_CONTEXT\_WINDOWS}, \code{COMPLETED}.
        Раздаёт работу N воркерам и аккумулирует ответы.
  \item \code{DslExecutionWorkerActor}: вычисляет один кандидат
        либо одно окно; для \code{DISTANCE} обсчитывает пары.
  \item \code{DslExecutionResultCollectorActor}: финальный
        результат.
\end{itemize}

\uml{dsl\_executor\_sequence.puml}

\subsection{Зачем акторы}

\begin{itemize}
  \item Параллельная обработка независимых кандидатов и пар без явной
        работы с потоками.
  \item Чёткое разделение ответственностей: координатор не считает,
        воркер не маршрутизирует.
  \item Доменные FSM-состояния служат живой документацией: по списку
        состояний понятно, что вообще делает актор.
\end{itemize}

\uml{fanout\_sequence.puml}

\subsection{End-to-end поток}

Полный сквозной путь от пользовательского ввода до структурированного
результата:

\begin{enumerate}
  \item текст запроса;
  \item поток токенов после лексера;
  \item Query AST после парсера;
  \item DocumentIndex поверх AstDocument;
  \item результат в виде JSON.
\end{enumerate}

\uml{dsl\_sequence.puml}\ \uml{tex\_dsl\_query\_sequence.puml}

% =============================================================
\section{Примеры запросов и фактический вывод}

Все запросы взяты из \code{examples.md}. Запуск производится скриптом
\code{run\_examples.py} над документом \code{GA\_1\_2025.tex} (учебный
конспект на русском). Команда:

\begin{Verbatim}[fontsize=\small,frame=single]
PYTHONPATH=src python run_examples.py GA_1_2025.tex
\end{Verbatim}

Все четыре запроса завершились без ошибок. Ниже для каждого приведён
текст запроса и фрагмент фактического вывода (полные выгрузки в
\code{examples\_output.txt}).

\subsection{Пример 1. FIND: определения о языках}

\textbf{Запрос:}
\begin{Verbatim}[fontsize=\small,frame=single,formatcom=\color{dslkw}]
FIND definition
WHERE metadata.index ~= "Язык"
RETURN text, count
\end{Verbatim}

\textbf{Семантика:} выбрать все узлы сущности \code{definition},
у которых метаданные содержат поле \code{index}, сопоставимое regex
\code{Язык}; вернуть текст узла и общее количество.

\textbf{Вывод} (фрагмент, всего \code{count = 19}):
\begin{Verbatim}[fontsize=\scriptsize,frame=single]
{
  "type": "find_query_execution_result",
  "count": 19,
  "returns": ["text", "count"],
  "items": [
    { "text": "\\odmkey{языком}{Язык (language)}" },
    { "text": "\\odmkey{Языком}{Язык!порождаемый грамматикой}" },
    { "text": "\\odmkey{разрешимым}{Язык!разрешимый}" },
    { "text": "\\odmkey{легко разрешимым}{Язык!легко разрешимый}" },
    { "text": "\\odmkey{полуразрешимым}{Язык!полуразрешимый}" },
    { "text": "\\odmkey{распознающим}{Язык!распознаваемый}" },
    { "text": "\\odmkey{линейным}{Язык!линейный}" },
    { "text": "\\odmkey{леволинейным}{Язык!леволинейный}" },
    { "text": "\\odmkey{праволинейным}{Язык!праволинейный}" },
    { "text": "\\odmkey{языки разметки}{Язык!разметки}" },
    { "text": "\\odmkey{SGML}{Язык!разметки!SGML}" },
    { "text": "\\odmkey{DTD}{Язык!разметки!DTD}" },
    { "text": "\\odmkey{HTML}{Язык!разметки!HTML}" },
    { "text": "\\odmkey{XML}{Язык!разметки!XML}" },
    { "text": "\\odmkey{XML Schema}{Язык!разметки!XML Schema}" },
    { "text": "\\odmkey{элементарным}{Язык!элементарный}" },
    { "text": "\\odmkey{регулярным}{Язык!регулярный}" },
    { "text": "\\odmkey{языка Дика}{Язык!Дика}" }
  ]
}
\end{Verbatim}

\subsection{Пример 2. DISTANCE между определениями (в символах)}

\textbf{Запрос:}
\begin{Verbatim}[fontsize=\small,frame=single,formatcom=\color{dslkw}]
DISTANCE definition [metadata.index ~= "Язык"]
TO       definition [metadata.index ~= "Язык"]
LIMIT_PAIRS all
RETURN pairs, stats, distance(symbol), count
\end{Verbatim}

\textbf{Семантика:} построить все пары определений со словом «Язык»
в индексе и измерить расстояние между ними в символах текста.
Вернуть список пар, агрегированную статистику и количество.

\textbf{Вывод} (фрагмент, всего 171 пара):
\begin{Verbatim}[fontsize=\scriptsize,frame=single]
{
  "type": "distance_query_execution_result",
  "count": 171,
  "returns": ["pairs", "stats", "distance(symbol)", "count"],
  "items": [
    {
      "left":  { "entity": "definition",
                 "text": "\\odmkey{леволинейным}{Язык!леволинейный}",
                 "start": 56900, "end": 56940 },
      "right": { "entity": "definition",
                 "text": "\\odmkey{праволинейным}{Язык!праволинейный}",
                 "start": 56945, "end": 56987 },
      "span":  { "start": 56900, "end": 56987 },
      "distance": { "unit": "symbol", "value": 5 }
    },
    ...
  ],
  "stats": {
    "count": 171,
    "mean": 35277.4,
    "variance": 7.75e8,
    "min": 5,
    "max": 136374
  }
}
\end{Verbatim}

\textbf{Интерпретация:} ближайшая пара отстоит на 5 символов
(соседние определения в одной строке), самая далёкая на 136 тысяч
символов; среднее расстояние около 35 тысяч символов.

\subsection{Пример 3. DISTANCE между определениями (в определениях)}

\textbf{Запрос:}
\begin{Verbatim}[fontsize=\small,frame=single,formatcom=\color{dslkw}]
DISTANCE definition [metadata.index ~= "Язык"]
TO       definition [metadata.index ~= "Язык"]
LIMIT_PAIRS all
RETURN pairs, stats, distance(definition), count
\end{Verbatim}

\textbf{Семантика:} те же пары, но единица измерения иная: количество
\code{definition}-узлов, лежащих строго между членами пары. Это даёт
структурную, а не текстовую дистанцию.

\textbf{Вывод} (фрагмент):
\begin{Verbatim}[fontsize=\scriptsize,frame=single]
{
  "type": "distance_query_execution_result",
  "count": 171,
  "returns": ["pairs", "stats", "distance(definition)", "count"],
  "items": [
    {
      "left":  { "entity": "definition",
                 "text": "\\odmkey{разрешимым}{Язык!разрешимый}",
                 "start": 42983, "end": 43019 },
      "right": { "entity": "definition",
                 "text": "\\odmkey{легко разрешимым}{Язык!легко разрешимый}",
                 "start": 43396, "end": 43444 },
      "span":  { "start": 42983, "end": 43444 },
      "distance": { "unit": "definition", "value": 0 }
    },
    ...
  ],
  "stats": {
    "count": 171,
    "mean": 35.0,
    "variance": 719.7,
    "min": 0,
    "max": 132
  }
}
\end{Verbatim}

\textbf{Интерпретация:} часть пар стоит подряд (между ними нет других
определений, distance = 0), максимум есть 132 промежуточных
определения; среднее около 35.

\subsection{Пример 4. DISTANCE: определение и теорема}

\textbf{Запрос:}
\begin{Verbatim}[fontsize=\small,frame=single,formatcom=\color{dslkw}]
DISTANCE definition    [metadata.index ~= "Язык"]
TO       semantic_block [metadata.kind = "theorem"]
LIMIT_PAIRS all
RETURN pairs, stats, distance(symbol), count
\end{Verbatim}

\textbf{Семантика:} измерить расстояние между каждым определением о
языке и каждой теоремой документа. Полезно для оценки, как близко
формальные утверждения стоят к вводимым понятиям.

\textbf{Вывод} (фрагмент, всего 798 пар):
\begin{Verbatim}[fontsize=\scriptsize,frame=single]
{
  "type": "distance_query_execution_result",
  "count": 798,
  "returns": ["pairs", "stats", "distance(symbol)", "count"],
  "items": [
    {
      "left":  { "entity": "definition",
                 "text": "\\odmkey{распознающим}{Язык!распознаваемый}",
                 "start": 43686, "end": 43728 },
      "right": { "entity": "semantic_block",
                 "text": "\\begin{theorem} Язык ... \\end{theorem}",
                 "start": 43732, "end": 43889 },
      "span":  { "start": 43686, "end": 43889 },
      "distance": { "unit": "symbol", "value": 1 }
    },
    ...
  ],
  "stats": {
    "count": 798,
    "mean": 89318.2,
    "variance": 3.88e9,
    "min": 1,
    "max": 268827
  }
}
\end{Verbatim}

\textbf{Интерпретация:} некоторые теоремы стоят буквально вплотную
к определениям (1 символ), но среднее расстояние велико (около
89 тысяч символов), поскольку каждый из 19 определений породил пару
с каждой теоремой документа (798 пар).

% =============================================================
\section{Заключение}

\subsection{Итоги}

\begin{itemize}
  \item Реализован декларативный язык запросов над AST документа с
        тремя типами запросов (\code{FIND}, \code{CONTEXT},
        \code{DISTANCE}) и богатыми \code{WHERE}- и
        \code{RETURN}-спецификациями.
  \item Архитектура построена строго снизу вверх: actor framework,
        затем document AST, затем DSL. Это исключает обратные
        зависимости и упрощает тестирование.
  \item Recursive-descent парсер и типобезопасный AST упростили
        добавление новых конструкций (как было с \code{DISTANCE}).
  \item Акторный конвейер с явными FSM делает поведение системы
        предсказуемым и допускает параллельную обработку независимых
        кандидатов.
  \item Все четыре примера из \code{examples.md} успешно отрабатывают
        на \code{GA\_1\_2025.tex}: 19 определений со словом «Язык»,
        171 пара определений (с двумя единицами измерения), 798 пар
        вида «определение, теорема».
\end{itemize}

\subsection{Возможные направления развития}

\begin{itemize}
  \item Поддержка дополнительных встроенных функций в \code{WHERE}
        (например, \code{count\_within}, \code{depth}).
  \item Кэширование DocumentIndex между запросами в интерактивном
        режиме CLI.
  \item Параллельный запуск конвейера исполнения с реальным пулом
        потоков (сейчас драйвер пошаговый).
  \item Расширение набора форматов (Markdown, XML) через новые
        конфиги в \code{config/formats/}.
\end{itemize}

\subsection{Указатель UML-диаграмм}

Все файлы лежат в \code{docs/uml/}:

\begin{itemize}
  \item \code{component\_overview.puml}: общая компонентная
        диаграмма системы.
  \item \code{dsl\_architecture.puml}: классы и слои DSL.
  \item \code{dsl\_parser\_sequence.puml}: последовательность
        разбора запроса.
  \item \code{dsl\_executor\_sequence.puml}: последовательность
        исполнения запроса.
  \item \code{dsl\_sequence.puml}: end-to-end поток от текста
        запроса до результата.
  \item \code{actor\_architecture.puml}: устройство акторного
        фреймворка.
  \item \code{fanout\_sequence.puml}: fan-out от координатора
        к воркерам.
  \item \code{tex\_dsl\_query\_sequence.puml}: конкретный
        сценарий на TeX-документе.
  \item \code{tex\_pipeline\_overview.puml},
        \code{tex\_parsing\_sequence.puml},
        \code{tex\_ast\_model.puml}: TeX-специфичная часть
        пайплайна.
  \item \code{ast\_parser\_architecture.puml},
        \code{ast\_parser\_sequence.puml}: акторный парсер
        документа.
\end{itemize}

\end{document}
